\documentclass[11pt]{article}

% Page setup
\usepackage[margin=1in]{geometry}
\usepackage{hyperref}
\hypersetup{
    colorlinks=true,
    linkcolor=blue,
    urlcolor=blue,
    citecolor=blue
}

% Graphics and code
\usepackage{graphicx}
\usepackage{listings}
\usepackage{xcolor}
\usepackage{float}

% Tables
\usepackage{booktabs}
\usepackage{longtable}
\usepackage{array}
\usepackage{colortbl}

% Fonts and formatting
\usepackage{enumitem}
\usepackage{parskip}

% Code listing style
\definecolor{codebg}{RGB}{248,248,248}
\definecolor{codegreen}{RGB}{0,128,0}
\definecolor{codegray}{RGB}{128,128,128}
\definecolor{codeblue}{RGB}{0,0,192}
\definecolor{codepurple}{RGB}{128,0,128}

\lstdefinelanguage{Go}{
    keywords={break, case, chan, const, continue, default, defer, else, fallthrough, for, func, go, goto, if, import, interface, map, package, range, return, select, struct, switch, type, var},
    keywordstyle=\color{codeblue}\bfseries,
    keywords=[2]{string, int, int64, bool, error, nil, true, false},
    keywordstyle=[2]\color{codepurple},
    identifierstyle=\color{black},
    sensitive=true,
    comment=[l]{//},
    morecomment=[s]{/*}{*/},
    commentstyle=\color{codegreen}\ttfamily,
    stringstyle=\color{codegray},
    morestring=[b]",
    morestring=[b]`
}

\lstset{
    language=Go,
    backgroundcolor=\color{codebg},
    basicstyle=\ttfamily\small,
    breakatwhitespace=false,
    breaklines=true,
    captionpos=b,
    keepspaces=true,
    numbers=left,
    numbersep=8pt,
    numberstyle=\tiny\color{codegray},
    showspaces=false,
    showstringspaces=false,
    showtabs=false,
    tabsize=4,
    frame=single,
    framesep=5pt,
    xleftmargin=15pt
}

\title{\textbf{ClusterShip} \\ \large A Terminal Battleship Game for Learning Kubernetes \\ \vspace{0.5cm} \normalsize Game Guide \& Code Walkthrough}
\author{ClusterShip Documentation}
\date{}

\begin{document}

\maketitle

\begin{abstract}
ClusterShip is a terminal-based battleship game that teaches Kubernetes concepts through interactive gameplay. This document serves as both a player-facing game guide and a developer-facing code walkthrough, explaining how game mechanics map to real Kubernetes architecture, the design decisions behind each package, and how the optional real-cluster integration works.
\end{abstract}

\tableofcontents
\newpage

\section{Introduction}

\subsection{What is ClusterShip?}

ClusterShip is a terminal battleship game where you battle AI opponents representing major tech companies (Netflix, AWS, Google, Meta, etc.) on a shared 100$\times$100 ocean. Each company operates a ``fleet'' of regions (data centers), racks (server slots), and pods (workload units) that mirror real Kubernetes infrastructure.

The game teaches Kubernetes concepts through gameplay:
\begin{itemize}
    \item \textbf{Pods} are individual workload units that can be destroyed and rescheduled
    \item \textbf{Services} maintain desired replica counts and handle failover
    \item \textbf{Affinity rules} determine whether pods can be rescheduled when their rack is destroyed
    \item \textbf{Namespaces} isolate each company's resources
    \item \textbf{Events} show real-time pod lifecycle changes
\end{itemize}

\subsection{Simulation vs Real Kubernetes}

ClusterShip operates in two modes:

\begin{enumerate}
    \item \textbf{Simulation Mode (default)}: All game logic runs locally. Pod scheduling, attacks, and rescheduling are simulated in memory. No actual Kubernetes cluster is required.

    \item \textbf{Real Kubernetes Mode}: When enabled in settings, ClusterShip deploys actual pods to your Kubernetes cluster. Attacks translate to \texttt{kubectl delete} commands, and pod events stream from the real cluster's API server.
\end{enumerate}

\begin{figure}[H]
    \centering
    \fbox{\parbox{0.8\textwidth}{\centering\vspace{2cm}\textbf{Figure: Architecture Overview}\\Game Layer $\rightarrow$ K8s Integration Layer $\rightarrow$ External Cluster\vspace{2cm}}}
    \caption{ClusterShip architecture showing game layer, optional K8s integration, and external cluster communication.}
    \label{fig:architecture}
\end{figure}

\subsection{High-Level Gameplay Loop}

\begin{enumerate}
    \item \textbf{Menu}: Choose New Game, Demo Mode, Settings, or Tutorial
    \item \textbf{Company Select}: Pick your company (e.g., Netflix, AWS)
    \item \textbf{Enemy Select}: Choose 1--5 AI opponents
    \item \textbf{Placement}: Fleets are placed on the shared ocean (no overlap)
    \item \textbf{Battle}: Take turns attacking enemy racks; AI opponents attack you
    \item \textbf{Game Over}: Win by destroying all enemy pods, or lose when yours are gone
\end{enumerate}

\section{Getting Started \& Configuration}

\subsection{Prerequisites}

\begin{itemize}
    \item \textbf{Go 1.24+}: Required to build and run
    \item \textbf{Terminal}: Any terminal supporting ANSI colors (iTerm2, Windows Terminal, etc.)
    \item \textbf{Optional}: Docker Desktop or \texttt{kind} for real K8s integration
\end{itemize}

\subsection{Installation \& Running}

\begin{lstlisting}[language=bash,numbers=none]
# Clone the repository
git clone https://github.com/ndg8743/ClusterShip.git
cd ClusterShip

# Run directly
go run ./cmd/clustership

# Or build and run
go build -o clustership ./cmd/clustership
./clustership
\end{lstlisting}

\subsection{Configuration}

Settings are stored in \texttt{\textasciitilde/.clustership/config.json}. The \texttt{GameConfig} struct in \texttt{pkg/config/config.go} defines all configurable options:

\begin{lstlisting}
type GameConfig struct {
    BoardWidth     int64  `json:"board_width"`
    BoardHeight    int64  `json:"board_height"`
    ShipsPerPlayer int    `json:"ships_per_player"`
    RacksPerShip   int    `json:"racks_per_ship"`
    PodsPerRack    int    `json:"pods_per_rack"`
    TurnDelayMs    int    `json:"turn_delay_ms"`
    EnableRealK8s  bool   `json:"enable_real_k8s"`
    K8sNamespace   string `json:"k8s_namespace"`
    Kubeconfig     string `json:"kubeconfig"`
    UseSparseBoard bool   `json:"use_sparse_board"`
}
\end{lstlisting}

Access these via the in-game Settings menu (\texttt{S} key from main menu).


\section{Game Concepts and Kubernetes Mapping}


Every ClusterShip game element has a direct Kubernetes equivalent. Understanding these mappings helps you learn real K8s concepts through gameplay.

\subsection{Component Mapping Table}

\begin{table}[H]
\centering
\begin{tabular}{@{}lll@{}}
\toprule
\textbf{ClusterShip} & \textbf{Kubernetes} & \textbf{Role} \\
\midrule
Board & API Server + etcd & Single source of truth for game state \\
Company & Namespace & Isolated group of resources \\
Region (Ship) & Node & Physical machine / data center \\
Rack (Cell) & Node capacity slot & Individual server slot \\
Pod & Pod & Workload unit that can be killed/rescheduled \\
Service & Deployment + Service & Logical grouping with replica count \\
AI Player & Controller/Operator & Watches state, makes decisions \\
Attack & \texttt{kubectl delete pod} & Destroys workloads \\
Rescheduling & Pod eviction/recreation & Moves pods based on affinity \\
\bottomrule
\end{tabular}
\caption{ClusterShip to Kubernetes component mapping}
\label{tab:mapping}
\end{table}

\subsection{Conceptual Hierarchy}

The game follows a strict hierarchy that mirrors K8s resource ownership:

\begin{verbatim}
Company (Namespace)
  |-- Service (Deployment)
  |     |-- Pod 1
  |     |-- Pod 2
  |     `-- Pod 3
  `-- Region (Node)
        |-- Rack 0 --> [Pod, Pod]
        |-- Rack 1 --> [Pod]
        `-- Rack 2 --> [Pod]
\end{verbatim}

% Placeholder for hierarchy diagram
\begin{figure}[H]
    \centering
    \fbox{\parbox{0.8\textwidth}{\centering\vspace{2cm}\textbf{Figure: Conceptual Hierarchy}\\Company $\rightarrow$ Services $\rightarrow$ Pods $\rightarrow$ Regions $\rightarrow$ Racks\vspace{2cm}}}
    \caption{Hierarchical relationship between game components}
    \label{fig:hierarchy}
\end{figure}

\subsection{Affinity Types}

Affinity determines what happens when a rack is destroyed:

\begin{table}[H]
\centering
\begin{tabular}{@{}llll@{}}
\toprule
\textbf{Affinity} & \textbf{Icon} & \textbf{Rescheduling} & \textbf{K8s Equivalent} \\
\midrule
Hard & {[!]} & NO - pod stays Pending & \texttt{requiredDuringScheduling} \\
Soft & {[o]} & Yes, prefers same region & \texttt{preferredDuringScheduling} \\
Spread & {[\textasciitilde]} & Yes, spreads across racks & \texttt{podAntiAffinity} \\
None & {[-]} & Yes, anywhere & No constraints \\
\bottomrule
\end{tabular}
\caption{Affinity types and their Kubernetes equivalents}
\label{tab:affinity}
\end{table}

Hard affinity is used for stateful workloads like primary databases that cannot be moved. Spread affinity is used for CDN edge nodes that benefit from distribution. This mirrors real-world K8s scheduling decisions.


\section{Companies, Services, and Templates}


\subsection{Company Templates}

Companies are defined in JSON template files under \texttt{templates/companies/}. Each template specifies the company's regions, services, AI strategy, and difficulty.

\begin{lstlisting}
type CompanyTemplate struct {
    ID          string           `json:"id"`
    Name        string           `json:"name"`
    Emoji       string           `json:"emoji"`
    Description string           `json:"description"`
    Regions     []RegionTemplate `json:"regions"`
    Services    []ServiceTemplate `json:"services"`
    AIStrategy  AIStrategy       `json:"ai_strategy"`
    Difficulty  string           `json:"difficulty"`
}
\end{lstlisting}

\subsection{Service Templates}

Services define logical groupings of pods with shared characteristics:

\begin{lstlisting}
type ServiceTemplate struct {
    ID             string       `json:"id"`
    Name           string       `json:"name"`
    Replicas       int          `json:"replicas"`
    PodsPerReplica int          `json:"pods_per_replica"`
    Affinity       AffinityType `json:"affinity"`
    Criticality    string       `json:"criticality"`
    CanFailover    bool         `json:"can_failover"`
}
\end{lstlisting}

\subsection{Example Companies and Services}

\begin{table}[H]
\centering
\begin{tabular}{@{}llll@{}}
\toprule
\textbf{Company} & \textbf{Service} & \textbf{Affinity} & \textbf{Description} \\
\midrule
Netflix & cdn-edge & Spread & CDN nodes distributed globally \\
Netflix & database & Hard & Primary database, cannot move \\
Netflix & playback & Soft & Video playback servers \\
AWS & ec2-api & Soft & Compute API servers \\
AWS & s3-storage & Hard & Object storage, stateful \\
Google & bigtable & Hard & Distributed database \\
\bottomrule
\end{tabular}
\caption{Example companies and their services}
\label{tab:companies}
\end{table}

\subsection{K8s Manifest Templates}

When real K8s mode is enabled, each service has corresponding YAML manifests in \texttt{templates/k8s/<company>/}. These are real Kubernetes Deployments and Services that run actual containers:

\begin{verbatim}
templates/k8s/
  netflix/
    cdn-edge.yaml      # nginx + curl sidecar
    database.yaml      # postgres
    playback.yaml      # node.js
  aws/
    ec2-api.yaml       # python HTTP
    s3-storage.yaml    # minio
\end{verbatim}


\section{Playing the Game: Controls and Views}


\subsection{Controls}

\begin{table}[H]
\centering
\begin{tabular}{@{}ll@{}}
\toprule
\textbf{Key} & \textbf{Action} \\
\midrule
Arrow keys / WASD & Move cursor on the board \\
Enter / Space & Fire at cursor position \\
1--5 & Change view level \\
C & Cycle service display (your services vs enemy) \\
D & Toggle debug mode (see all ships) \\
I & Toggle info overlay \\
Q & Quit / Back to menu \\
\bottomrule
\end{tabular}
\caption{Game controls}
\label{tab:controls}
\end{table}

\subsection{View Levels}

The game provides five hierarchical view levels, accessed with number keys 1--5:

\begin{enumerate}
    \item \textbf{Map View (1)}: Overview of the 100$\times$100 ocean showing ships and attacks
    \item \textbf{Ship View (2)}: List of your regions (ships) with health status
    \item \textbf{Rack View (3)}: Drill into individual racks within a ship, showing pods
    \item \textbf{YAML View (4)}: K8s manifests (fog of war applied to enemies)
    \item \textbf{Rack Layout (5)}: Visual grid showing pod distribution per rack
\end{enumerate}

\subsection{Game Interface}

\begin{verbatim}
+-------------------------------------------------------------+
| CLUSTERSHIP - Player vs Netflix vs AWS    Turn: 42 [VIEW 1] |
+-------------------------------------------------------------+
|                                                              |
| 100x100 SHARED OCEAN        SERVICE STATUS                   |
| +-------------------+       [!]Origin API  [####-] 80%       |
| | . . . . X . . . . |       [~]CDN Edge    [#####] 100%      |
| | . . # # # # # . . |       [o]Playback    [###--] 60%       |
| | . . . . . . . . . |                                        |
| | . . . . . O . . . | <cursor                                |
| +-------------------+       K8S EVENTS                       |
|                             [+] Pod cdn-edge rescheduled     |
| [1-5] view  [arrows] move  [enter] fire  [q] quit           |
+-------------------------------------------------------------+
\end{verbatim}


\section{Code Architecture Overview}


\subsection{Package Structure}

ClusterShip is organized into focused packages, each with a specific responsibility:

\begin{verbatim}
cmd/clustership/     Entry point (main.go)
pkg/
  config/            Configuration loading, validation, persistence
  game/              Domain model: Company, Region, Rack, Pod, Service
  sparse/            Sparse board representation using quadtree
  tui/               Terminal UI using Bubble Tea framework
  k8s/               Real Kubernetes client, deployer, watcher
  benchmark/         Performance benchmarking and metrics
  hardware/          System detection for performance tiers
\end{verbatim}

% Placeholder for package dependency diagram
\begin{figure}[H]
    \centering
    \fbox{\parbox{0.8\textwidth}{\centering\vspace{2cm}\textbf{Figure: Package Dependency Graph}\\cmd $\rightarrow$ tui $\rightarrow$ game, config, k8s, benchmark\\tui $\rightarrow$ sparse\vspace{2cm}}}
    \caption{Package dependencies in ClusterShip}
    \label{fig:packages}
\end{figure}

\subsection{Why These Packages?}

\begin{itemize}
    \item \textbf{config}: Centralized configuration with hardware-aware defaults. Uses Go's \texttt{encoding/json} for persistence.

    \item \textbf{game}: Pure domain model with no UI or I/O dependencies. Makes testing trivial and allows reuse.

    \item \textbf{sparse}: For large boards (1000$\times$1000+), storing every cell is wasteful. The quadtree provides O(log n) lookups with minimal memory.

    \item \textbf{tui}: Bubble Tea is chosen for its functional architecture (Model-Update-View) and excellent terminal rendering via Lip Gloss.

    \item \textbf{k8s}: Isolates all Kubernetes dependencies. Uses \texttt{client-go} for API communication, which is the official Go client library.

    \item \textbf{benchmark}: Supports stress testing large boards and measuring performance. Uses atomic counters for thread-safe metrics.

    \item \textbf{hardware}: Detects RAM, CPU cores, and GPU availability to set appropriate defaults for the user's system.
\end{itemize}

\subsection{Entry Point}

\begin{lstlisting}
// cmd/clustership/main.go
func main() {
    p := tea.NewProgram(tui.NewAppModel(), tea.WithAltScreen())
    if _, err := p.Run(); err != nil {
        fmt.Printf("Error: %v\n", err)
        os.Exit(1)
    }
}
\end{lstlisting}

The entry point is minimal---it creates a Bubble Tea program with the \texttt{AppModel} and runs it. All game logic lives in the \texttt{tui} package.


\section{The Game Model: Companies, Regions, Racks, Pods}


\subsection{Core Types}

The \texttt{pkg/game/company.go} file defines the core domain types:

\begin{lstlisting}
type Company struct {
    ID          string
    Name        string
    Regions     []*Region
    Services    []*Service
    AIStrategy  AIStrategy
}

type Region struct {
    ID        string
    Name      string
    RackCount int
    Racks     []*Rack
    Placement [][2]int  // board coordinates
    IsDestroyed bool
}

type Rack struct {
    ID          string
    Position    [2]int
    Capacity    int
    Pods        []*Pod
    IsDestroyed bool
}

type Pod struct {
    ID        string
    ServiceID string
    Health    int
    Status    PodStatus  // Running, Pending, Terminated
}
\end{lstlisting}

\subsection{Services and Affinity}

Services represent logical workload groupings:

\begin{lstlisting}
type Service struct {
    ID          string
    Name        string
    Replicas    int
    Affinity    AffinityType
    Pods        []*Pod
    IsHealthy   bool
}

type AffinityType string

const (
    AffinityHard   AffinityType = "hard"
    AffinitySoft   AffinityType = "soft"
    AffinitySpread AffinityType = "spread"
    AffinityNone   AffinityType = "none"
)
\end{lstlisting}

\subsection{Attack Flow}

When an attack occurs:
\begin{enumerate}
    \item Board identifies the cell at (x, y)
    \item If cell has a rack, damage is applied to pods on that rack
    \item If pod health reaches 0, check service affinity
    \item Hard affinity: pod stays Pending (cannot reschedule)
    \item Other affinity: find new rack and reschedule pod
    \item Generate GameEvent for the event log
\end{enumerate}


\section{The Board and Sparse Representation}


\subsection{Why Sparse?}

A 100$\times$100 board has 10,000 cells, which is manageable. But ClusterShip supports boards up to 10,000,000$\times$10,000,000 for stress testing. A dense array would require terabytes of memory. Instead, we use a quadtree.

\subsection{Quadtree Implementation}

The \texttt{pkg/sparse/quadtree.go} implements a region quadtree:

\begin{lstlisting}
type QuadTree struct {
    root   *node
    bounds Bounds
    size   int
}

type node struct {
    bounds   Bounds
    children [4]*node  // NW, NE, SW, SE
    items    []Item
    isLeaf   bool
}

func (qt *QuadTree) Insert(x, y int64, value interface{}) {
    // Recursively subdivide until item fits in a small enough region
}

func (qt *QuadTree) Query(x, y int64) (interface{}, bool) {
    // O(log n) lookup
}
\end{lstlisting}

\subsection{Board Interface}

The \texttt{pkg/sparse/board.go} provides the game-facing interface:

\begin{lstlisting}
type SparseBoard struct {
    tree   *QuadTree
    width  int64
    height int64
}

func (b *SparseBoard) SetCell(x, y int64, cell *Cell) error
func (b *SparseBoard) GetCell(x, y int64) (*Cell, bool)
func (b *SparseBoard) CellCount() int
\end{lstlisting}

% Placeholder for quadtree diagram
\begin{figure}[H]
    \centering
    \fbox{\parbox{0.8\textwidth}{\centering\vspace{2cm}\textbf{Figure: Quadtree Visualization}\\2$\times$2 recursive subdivision showing occupied vs empty cells\vspace{2cm}}}
    \caption{Quadtree spatial subdivision}
    \label{fig:quadtree}
\end{figure}


\section{TUI and Game Loop}


\subsection{Bubble Tea Architecture}

ClusterShip uses the \href{https://github.com/charmbracelet/bubbletea}{Bubble Tea} framework, which implements the Elm Architecture:

\begin{itemize}
    \item \textbf{Model}: Application state (\texttt{AppModel})
    \item \textbf{Update}: Handle messages, return new model + commands
    \item \textbf{View}: Render model to string for display
\end{itemize}

\begin{lstlisting}
type AppModel struct {
    state        GameState
    cfg          *config.GameConfig
    board        *Board
    ais          map[string]*AIPlayer
    cursor       [2]int
    turn         int
    k8sClient    *k8s.Client
    // ... more fields
}

func (m AppModel) Init() tea.Cmd {
    return doTick()
}

func (m AppModel) Update(msg tea.Msg) (tea.Model, tea.Cmd) {
    switch msg := msg.(type) {
    case tea.KeyMsg:
        return m.handleKeyPress(msg)
    case tickMsg:
        return m.handleTick()
    }
    return m, nil
}

func (m AppModel) View() string {
    return m.renderCurrentState()
}
\end{lstlisting}

\subsection{Game States}

The game uses a finite state machine:

\begin{lstlisting}
type GameState int

const (
    StateMenu GameState = iota
    StateCompanySelect
    StateEnemySelect
    StatePlacement
    StateBattle
    StateGameOver
    StateSettings
    // ...
)
\end{lstlisting}

% Placeholder for state machine diagram
\begin{figure}[H]
    \centering
    \fbox{\parbox{0.8\textwidth}{\centering\vspace{2cm}\textbf{Figure: State Machine}\\Menu $\rightarrow$ CompanySelect $\rightarrow$ EnemySelect $\rightarrow$ Placement $\rightarrow$ Battle $\rightarrow$ GameOver\vspace{2cm}}}
    \caption{Game state transitions}
    \label{fig:states}
\end{figure}

\subsection{Tick System}

The game loop is driven by ticks:

\begin{lstlisting}
type tickMsg time.Time

func doTickWithDelay(ms int) tea.Cmd {
    return tea.Tick(time.Duration(ms)*time.Millisecond,
        func(t time.Time) tea.Msg {
            return tickMsg(t)
        })
}

func (m *AppModel) tick() tea.Cmd {
    return doTickWithDelay(m.cfg.TurnDelayMs)
}
\end{lstlisting}

Each tick advances the turn counter and triggers AI moves during battle.


\section{Real Kubernetes Integration}


\subsection{When K8s Integration is Enabled}

In Settings, toggle ``Real K8s'' to ON. The game will:
\begin{enumerate}
    \item Create a namespace (\texttt{clustership} by default)
    \item Deploy all company manifests as real pods
    \item Watch pod events from the cluster
    \item Translate attacks to \texttt{kubectl delete} operations
    \item Clean up namespace on game exit
\end{enumerate}

\subsection{K8s Client}

The \texttt{pkg/k8s/client.go} wraps the official \texttt{client-go} library:

\begin{lstlisting}
type Client struct {
    clientset kubernetes.Interface
    config    *rest.Config
    namespace string
}

func NewClient(kubeconfig, namespace string) (*Client, error) {
    config, err := clientcmd.BuildConfigFromFlags("", kubeconfig)
    if err != nil {
        return nil, err
    }

    clientset, err := kubernetes.NewForConfig(config)
    if err != nil {
        return nil, err
    }

    return &Client{clientset: clientset, config: config, namespace: namespace}, nil
}
\end{lstlisting}

\subsection{Why client-go?}

\texttt{client-go} is the official Kubernetes Go client maintained by the Kubernetes project. It provides:
\begin{itemize}
    \item Type-safe API for all K8s resources
    \item Automatic authentication via kubeconfig
    \item Watch support for streaming events
    \item Retry and backoff handling
\end{itemize}

\subsection{Deployer}

The \texttt{pkg/k8s/deployer.go} handles resource creation:

\begin{lstlisting}
func (c *Client) DeployManifest(ctx context.Context, manifest *Manifest) error {
    switch manifest.Kind {
    case "Deployment":
        return c.deployDeployment(ctx, manifest)
    case "StatefulSet":
        return c.deployStatefulSet(ctx, manifest)
    case "Pod":
        return c.deployPod(ctx, manifest)
    case "Service":
        return c.deployService(ctx, manifest)
    }
    return nil
}
\end{lstlisting}

\subsection{Pod Watcher}

Real-time pod events are streamed via the K8s watch API:

\begin{lstlisting}
func (c *Client) WatchPods(ctx context.Context, ns string) (<-chan PodEvent, error) {
    watcher, err := c.clientset.CoreV1().Pods(ns).Watch(ctx, metav1.ListOptions{
        LabelSelector: "app=clustership",
    })
    if err != nil {
        return nil, err
    }

    events := make(chan PodEvent)
    go func() {
        for event := range watcher.ResultChan() {
            events <- convertEvent(event)
        }
    }()
    return events, nil
}
\end{lstlisting}

% Placeholder for K8s integration flow diagram
\begin{figure}[H]
    \centering
    \fbox{\parbox{0.8\textwidth}{\centering\vspace{2cm}\textbf{Figure: K8s Integration Flow}\\TUI $\rightarrow$ K8s Client $\rightarrow$ API Server $\rightarrow$ Pods $\rightarrow$ Watcher $\rightarrow$ Events $\rightarrow$ TUI\vspace{2cm}}}
    \caption{Real Kubernetes integration flow}
    \label{fig:k8sflow}
\end{figure}


\section{AI, Affinity, and Rescheduling Mechanics}


\subsection{AI Strategies}

Each AI company uses a strategy defined in \texttt{pkg/tui/ai.go}:

\begin{lstlisting}
type AIStrategy string

const (
    AIRandom     AIStrategy = "random"     // Pure random targeting
    AIHunter     AIStrategy = "hunter"     // Hunt neighbors after hit
    AIDefensive  AIStrategy = "defensive"  // Protect critical services
    AIAggressive AIStrategy = "aggressive" // Target high-value first
)
\end{lstlisting}

\textbf{AIRandom}: Picks a random cell. Simple but unpredictable.

\textbf{AIHunter}: After scoring a hit, searches adjacent cells to finish off the ship. Mimics human strategy.

\textbf{AIAggressive}: Uses KNN (K-Nearest Neighbors) analysis to identify clusters of ships and targets high-probability areas.

\textbf{AIDefensive}: Prioritizes attacks on enemies threatening its critical services.

\subsection{Rescheduling Algorithm}

When a pod is killed:

\begin{enumerate}
    \item Check the service's \texttt{AffinityType}
    \item \textbf{Hard}: Mark pod as Pending, do not reschedule
    \item \textbf{Soft}: Search same region first, then other regions
    \item \textbf{Spread}: Find the rack with fewest pods
    \item \textbf{None}: Find any rack with capacity
    \item If a suitable rack is found, move the pod there
    \item If no capacity exists, mark pod as Pending
\end{enumerate}

\begin{lstlisting}
func (b *Board) reschedulePod(pod *Pod, svc *Service) bool {
    switch svc.Affinity {
    case AffinityHard:
        return false  // Cannot reschedule
    case AffinitySoft:
        return b.findRackInRegion(pod.RegionID) != nil
    case AffinitySpread:
        return b.findLeastLoadedRack() != nil
    case AffinityNone:
        return b.findAnyRack() != nil
    }
    return false
}
\end{lstlisting}


\section{Configuration, Performance, and Benchmarks}


\subsection{Hardware Detection}

The \texttt{pkg/hardware/detect.go} analyzes the system:

\begin{lstlisting}
type SystemInfo struct {
    TotalRAM   uint64
    CPUCores   int
    HasGPU     bool
    GPUMemory  uint64
}

type PerformanceTier string

const (
    TierBasic    PerformanceTier = "basic"
    TierStandard PerformanceTier = "standard"
    TierAdvanced PerformanceTier = "advanced"
    TierExtreme  PerformanceTier = "extreme"
)
\end{lstlisting}

Each tier has different limits for board size, ships, and pods to ensure smooth gameplay.

\subsection{Performance Tiers}

\begin{table}[H]
\centering
\begin{tabular}{@{}llll@{}}
\toprule
\textbf{Tier} & \textbf{RAM} & \textbf{Max Board} & \textbf{Sparse Mode} \\
\midrule
Basic & $<$4GB & 100$\times$100 & Auto at $>$1000 \\
Standard & 4--8GB & 1,000$\times$1,000 & Auto at $>$1000 \\
Advanced & 8--16GB & 100,000$\times$100,000 & Auto at $>$1000 \\
Extreme & $>$16GB & 10M$\times$10M & Auto at $>$1000 \\
\bottomrule
\end{tabular}
\caption{Performance tiers and their limits}
\label{tab:tiers}
\end{table}

\subsection{Benchmark Package}

The \texttt{pkg/benchmark} package provides stress testing:

\begin{lstlisting}
type Metrics struct {
    TotalOps    atomic.Int64
    WorkerCount atomic.Int32
    OpsPerSec   atomic.Int64
}

type Runner struct {
    workers map[string]*Worker
    metrics *Metrics
}
\end{lstlisting}

Atomic counters ensure thread-safe metrics collection across multiple goroutines.


\section{Benchmark Results}

This section presents performance benchmarks measuring ClusterShip's scaling characteristics across all performance tiers on different hardware configurations.

\subsection{Test Environment 1: High-End Desktop}

\begin{table}[H]
\centering
\renewcommand{\arraystretch}{1.3}
\begin{tabular}{|l|l|}
\hline
\multicolumn{2}{|c|}{\textbf{Desktop Hardware Configuration}} \\
\hline
\textbf{Component} & \textbf{Specification} \\
\hline
CPU & 32 Cores (x86\_64) \\
RAM & 196 GB DDR5 \\
GPU & NVIDIA GeForce RTX 5070 Ti \\
VRAM & 16,303 MB GDDR7 \\
\hline
\multicolumn{2}{|c|}{\textbf{Software Environment}} \\
\hline
OS & Windows 11 \\
Go Version & 1.24.11 \\
GOMAXPROCS & 32 \\
Detected Tier & Extreme \\
\hline
\end{tabular}
\caption{Desktop benchmark environment specifications}
\label{tab:testenv-desktop}
\end{table}

\begin{table}[H]
\centering
\renewcommand{\arraystretch}{1.4}
\begin{tabular}{|c|c|c|c|c|c|}
\hline
\rowcolor{codebg}
\textbf{Tier} & \textbf{Board Size} & \textbf{Creation} & \textbf{Memory} & \textbf{Lookups/sec} & \textbf{Status} \\
\hline
\cellcolor{green!20}Basic & 100 $\times$ 100 & $<$0.01 ms & 0.02 MB & 2.5M & \textcolor{green!60!black}{PASS} \\
\hline
\cellcolor{blue!20}Standard & 1,000 $\times$ 1,000 & 1.00 ms & 0.22 MB & 19.9M & \textcolor{green!60!black}{PASS} \\
\hline
\cellcolor{orange!20}Advanced & 10,000 $\times$ 10,000 & 9.62 ms & 2.03 MB & 19.9M & \textcolor{green!60!black}{PASS} \\
\hline
\cellcolor{red!20}Extreme & 100,000 $\times$ 100,000 & 65.29 ms & 19.19 MB & 19.7M & \textcolor{green!60!black}{PASS} \\
\hline
\end{tabular}
\caption{Desktop performance benchmarks across all tiers}
\label{tab:benchmarks-desktop}
\end{table}

\subsection{Test Environment 2: Apple Silicon Laptop}

\begin{table}[H]
\centering
\renewcommand{\arraystretch}{1.3}
\begin{tabular}{|l|l|}
\hline
\multicolumn{2}{|c|}{\textbf{Laptop Hardware Configuration}} \\
\hline
\textbf{Component} & \textbf{Specification} \\
\hline
CPU & Apple M1 Pro (10 Cores) \\
RAM & 16 GB Unified Memory \\
GPU & Apple M1 Pro (16-core GPU) \\
\hline
\multicolumn{2}{|c|}{\textbf{Software Environment}} \\
\hline
OS & macOS \\
Go Version & 1.24.11 \\
GOMAXPROCS & 10 \\
Detected Tier & Standard \\
\hline
\end{tabular}
\caption{Laptop benchmark environment specifications}
\label{tab:testenv-laptop}
\end{table}

\begin{table}[H]
\centering
\renewcommand{\arraystretch}{1.4}
\begin{tabular}{|c|c|c|c|c|c|}
\hline
\rowcolor{codebg}
\textbf{Tier} & \textbf{Board Size} & \textbf{Creation} & \textbf{Memory} & \textbf{Lookups/sec} & \textbf{Status} \\
\hline
\cellcolor{green!20}Basic & 100 $\times$ 100 & --- & --- & --- & --- \\
\hline
\cellcolor{blue!20}Standard & 1,000 $\times$ 1,000 & --- & --- & --- & --- \\
\hline
\cellcolor{orange!20}Advanced & 10,000 $\times$ 10,000 & --- & --- & --- & --- \\
\hline
\cellcolor{red!20}Extreme & 100,000 $\times$ 100,000 & --- & --- & --- & --- \\
\hline
\end{tabular}
\caption{Laptop performance benchmarks across all tiers}
\label{tab:benchmarks-laptop}
\end{table}

\subsection{Scaling Analysis}

\begin{figure}[H]
\centering
\begin{tabular}{cc}
\fbox{\parbox{0.45\textwidth}{
\centering
\textbf{Board Creation Time (Desktop)}\\[0.5em]
\begin{tabular}{r|l}
Basic & \rule{0.1cm}{0.3cm} $<$0.01ms \\
Standard & \rule{0.5cm}{0.3cm} 1.00ms \\
Advanced & \rule{1.0cm}{0.3cm} 9.62ms \\
Extreme & \rule{2.5cm}{0.3cm} 65.29ms \\
\end{tabular}
}} &
\fbox{\parbox{0.45\textwidth}{
\centering
\textbf{Memory Usage (Desktop)}\\[0.5em]
\begin{tabular}{r|l}
Basic & \rule{0.1cm}{0.3cm} 0.02MB \\
Standard & \rule{0.2cm}{0.3cm} 0.22MB \\
Advanced & \rule{0.5cm}{0.3cm} 2.03MB \\
Extreme & \rule{2.5cm}{0.3cm} 19.19MB \\
\end{tabular}
}}
\end{tabular}
\caption{Visual comparison of creation time and memory usage across tiers (Desktop)}
\label{fig:benchviz}
\end{figure}

\subsection{Key Findings}

\begin{enumerate}
    \item \textbf{Efficient Sparse Representation}: The QuadTree-based sparse board enables a 100,000 $\times$ 100,000 board (10 billion potential cells) to use only 19.19 MB of memory.

    \item \textbf{Consistent Lookup Performance}: Attack lookups maintain 19--20 million operations per second across all tiers on desktop hardware, demonstrating O(log n) QuadTree query performance.

    \item \textbf{Linear Creation Scaling}: Board creation time scales linearly with the number of placed cells, not total board area:
    \begin{itemize}
        \item Basic: 300 cells $\rightarrow$ $<$0.01ms
        \item Standard: 1,000 cells $\rightarrow$ 1.00ms
        \item Advanced: 10,000 cells $\rightarrow$ 9.62ms
        \item Extreme: 100,000 cells $\rightarrow$ 65.29ms
    \end{itemize}

    \item \textbf{Memory Efficiency}: Memory grows with placed cell count, not board dimensions. The Extreme tier board is 1,000,000$\times$ larger in area than Basic but only uses approximately 960$\times$ more memory.
\end{enumerate}

\subsection{Running Your Own Benchmarks}

To run benchmarks on your system:

\begin{lstlisting}[language=bash,numbers=none]
# Build the benchmark tool
go build -o benchmark ./cmd/benchmark

# Run all tier benchmarks
./benchmark --all-tiers

# Run specific tier
./benchmark --tier Advanced

# Output as JSON for processing
./benchmark --all-tiers --json

# Or use Docker for isolated testing
docker build -f Dockerfile.benchmark -t clustership-benchmark .
docker run --rm --gpus all clustership-benchmark --all-tiers
\end{lstlisting}

\subsection{Performance Recommendations}

\begin{table}[H]
\centering
\begin{tabular}{|l|l|l|}
\hline
\textbf{System RAM} & \textbf{Recommended Tier} & \textbf{Max Board Size} \\
\hline
$<$ 4 GB & Basic & 100 $\times$ 100 \\
4--8 GB & Standard & 1,000 $\times$ 1,000 \\
8--16 GB & Advanced & 100,000 $\times$ 100,000 \\
$>$ 16 GB + GPU & Extreme & 10,000,000 $\times$ 10,000,000 \\
\hline
\end{tabular}
\caption{Recommended performance tier based on system resources}
\label{tab:recommendations}
\end{table}


\section{Extending ClusterShip}


\subsection{Adding a New Company}

\begin{enumerate}
    \item Create \texttt{templates/companies/mycompany.json}:
\begin{lstlisting}[language=json,numbers=none]
{
    "id": "mycompany",
    "name": "My Company",
    "regions": [{"id": "dc1", "name": "DC 1", "racks": 3}],
    "services": [{"id": "api", "name": "API", "replicas": 3, "affinity": "soft"}],
    "ai_strategy": "hunter"
}
\end{lstlisting}
    \item Create \texttt{templates/k8s/mycompany/api.yaml} (for real K8s mode)
    \item The game auto-discovers new companies on startup
\end{enumerate}

\subsection{Adding a New View}

\begin{enumerate}
    \item Add a new \texttt{ViewLevel} constant in \texttt{pkg/tui/app.go}
    \item Create a render function: \texttt{func (m *AppModel) renderMyView() string}
    \item Add case to the view switch in the main \texttt{View()} method
    \item Handle the key press to switch to your view
\end{enumerate}

\subsection{Custom Affinity Rules}

\begin{enumerate}
    \item Add new \texttt{AffinityType} constant in \texttt{pkg/game/company.go}
    \item Implement rescheduling logic in \texttt{pkg/tui/board.go}
    \item Update JSON schema to accept new affinity type
\end{enumerate}


\section{Dependencies and Libraries}


\subsection{Core Dependencies}

\begin{table}[H]
\centering
\begin{tabular}{@{}lll@{}}
\toprule
\textbf{Package} & \textbf{Purpose} & \textbf{Why Chosen} \\
\midrule
Bubble Tea & Terminal UI & Elm Architecture, functional, testable \\
Lip Gloss & Styling & Clean API for colors, borders, layout \\
client-go & K8s client & Official, type-safe, maintained by K8s \\
\bottomrule
\end{tabular}
\caption{Core dependencies}
\label{tab:deps}
\end{table}

\subsection{Bubble Tea Explained}

Bubble Tea implements the Elm Architecture in Go:
\begin{itemize}
    \item \textbf{Immutable state}: Each Update returns a new model
    \item \textbf{Message-driven}: All changes come through messages
    \item \textbf{Testable}: Pure functions make unit testing trivial
    \item \textbf{Composable}: Build complex UIs from simple components
\end{itemize}

\subsection{Lip Gloss Explained}

Lip Gloss provides declarative styling:
\begin{lstlisting}
var titleStyle = lipgloss.NewStyle().
    Bold(true).
    Foreground(lipgloss.Color("205")).
    Padding(0, 1)

rendered := titleStyle.Render("CLUSTERSHIP")
\end{lstlisting}


\section*{Appendix A: Reference Tables}

\addcontentsline{toc}{section}{Appendix A: Reference Tables}

\subsection*{All Companies}

\begin{longtable}{@{}lll@{}}
\toprule
\textbf{ID} & \textbf{Name} & \textbf{Difficulty} \\
\midrule
\endhead
netflix & Netflix & Medium \\
aws & Amazon Web Services & Hard \\
google & Google Cloud & Hard \\
meta & Meta & Medium \\
microsoft & Microsoft Azure & Hard \\
apple & Apple & Easy \\
\bottomrule
\end{longtable}

\subsection*{Key Types by Package}

\begin{longtable}{@{}ll@{}}
\toprule
\textbf{Package} & \textbf{Key Types} \\
\midrule
\endhead
pkg/game & Company, Region, Rack, Pod, Service, AffinityType, AIStrategy \\
pkg/tui & AppModel, GameState, ViewLevel, Board, AIPlayer \\
pkg/k8s & Client, Manifest, PodInfo, PodEvent \\
pkg/config & GameConfig \\
pkg/sparse & QuadTree, SparseBoard, Cell \\
pkg/benchmark & Metrics, Runner, Worker \\
pkg/hardware & SystemInfo, PerformanceTier, TierLimits \\
\bottomrule
\end{longtable}


\section*{References and Citations}

\addcontentsline{toc}{section}{References and Citations}

\begin{thebibliography}{99}




\bibitem{k8suprunning}
Burns, B., Beda, J., Hightower, K., \& Evenson, L. (2022). \textit{Kubernetes: Up and Running}, 3rd Edition.
O'Reilly Media. ISBN: 978-1098110208

\bibitem{dataintensive}
Kleppmann, M. (2017). \textit{Designing Data-Intensive Applications: The Big Ideas Behind Reliable, Scalable, and Maintainable Systems}.
O'Reilly Media. ISBN: 978-1449373320




\bibitem{schedulingsurvey}
Carrion, C. (2022). Kubernetes Scheduling: Taxonomy, Ongoing Issues and Challenges.
\textit{arXiv preprint arXiv:2203.00671}. \url{https://arxiv.org/abs/2203.00671}

\bibitem{k8sschedulingpaper}
Kubernetes Scheduling Algorithms Survey. (2025).
\textit{arXiv preprint arXiv:2504.17033}. \url{https://arxiv.org/pdf/2504.17033}

\bibitem{mincostflow}
Wikipedia. \textit{Minimum-cost flow problem}.
\url{https://en.wikipedia.org/wiki/Minimum-cost_flow_problem}

\bibitem{fibonacciheap}
Wikipedia. \textit{Fibonacci heap}.
\url{https://en.wikipedia.org/wiki/Fibonacci_heap}

\bibitem{dijkstra}
Modern Data 101. (2024). \textit{The New Dijkstra's Algorithm: Shortest Path Finding}.
\url{https://moderndata101.substack.com/p/the-new-dijkstras-algorithm-shortest}




\bibitem{bubbletea}
Charm. (2020--2024). \textit{Bubble Tea: A powerful little TUI framework for Go}.
GitHub. \url{https://github.com/charmbracelet/bubbletea}

\bibitem{lipgloss}
Charm. (2021--2024). \textit{Lip Gloss: Style definitions for nice terminal layouts in Go}.
GitHub. \url{https://github.com/charmbracelet/lipgloss}

\bibitem{gorillamux}
Gorilla Web Toolkit. \textit{gorilla/mux: A powerful HTTP router and URL matcher for Go}.
GitHub. \url{https://github.com/gorilla/mux}

\bibitem{clientgo}
Kubernetes Authors. (2016--2024). \textit{client-go: Go client for Kubernetes}.
GitHub. \url{https://github.com/kubernetes/client-go}




\bibitem{k8sdocs}
Kubernetes Authors. (2024). \textit{Kubernetes Documentation}.
\url{https://kubernetes.io/docs/home/}

\bibitem{gobyexample}
Warden, M. \textit{Go by Example}.
\url{https://gobyexample.com/}

\bibitem{minikubetutorials}
Kubernetes SIGs. \textit{minikube Tutorials}.
\url{https://minikube.sigs.k8s.io/docs/tutorials/}

\bibitem{golangdocs}
Go Authors. (2009--2024). \textit{The Go Programming Language}.
\url{https://go.dev}




\bibitem{k8sarchitecture}
Aqua Security. (2020). \textit{Kubernetes 101: Kubernetes Architecture}.
\url{https://www.aquasec.com/cloud-native-academy/kubernetes-101/kubernetes-architecture/}

\bibitem{k8scoreconcepts}
DZero Labs. \textit{Just-in-Time Kubernetes: A Beginner's Guide to Kubernetes Core Concepts}.
Medium. \url{https://medium.com/dzerolabs/just-in-time-kubernetes-a-beginners-guide-to-kubernetes-core-concepts-19ee7acbafa1}

\bibitem{k8sgoarchitecture}
LeapCell. \textit{Learning Large-Scale Go Project Architecture from K8s}.
Medium. \url{https://leapcell.medium.com/learning-large-scale-go-project-architecture-from-k8s-6c8f2c3862d8}

\bibitem{whatismux}
Kim, J. (2022). \textit{What Even Is A Mux?}
Dev.to. \url{https://dev.to/jpoly1219/what-even-is-a-mux-4fng}

\bibitem{containerdvsdocker}
Docker. \textit{containerd vs Docker}.
Docker Blog. \url{https://www.docker.com/blog/containerd-vs-docker/}

\bibitem{k8swiki}
Wikipedia. \textit{Kubernetes}.
\url{https://en.wikipedia.org/wiki/Kubernetes}




\bibitem{goplaylist}
Go Programming Tutorials. \textit{Go Programming Language Tutorial}.
YouTube Playlist. \url{https://youtube.com/playlist?list=PLmD8u-IFdreyh6EUfevBcbiuCKzFk0EW_}




\bibitem{quadtree}
Finkel, R. A., \& Bentley, J. L. (1974). Quad Trees: A Data Structure for Retrieval on Composite Keys.
\textit{Acta Informatica}, 4(1), 1--9. \url{https://doi.org/10.1007/BF00288933}

\bibitem{samet1990}
Samet, H. (1990). \textit{The Design and Analysis of Spatial Data Structures}.
Addison-Wesley. ISBN: 978-0201502558




\bibitem{elm}
Czaplicki, E. (2012--2024). \textit{The Elm Architecture}.
\url{https://guide.elm-lang.org/architecture/}

\bibitem{controllerpattern}
Kubernetes Authors. (2024). \textit{Kubernetes Controller Pattern}.
\url{https://kubernetes.io/docs/concepts/architecture/controller/}




\bibitem{k9s}
Poitras, F. (2019--2024). \textit{K9s: Kubernetes CLI To Manage Your Clusters In Style}.
GitHub. \url{https://github.com/derailed/k9s}

\bibitem{lazydocker}
Jesseduffield. (2019--2024). \textit{Lazydocker: A simple terminal UI for Docker and Docker Compose}.
GitHub. \url{https://github.com/jesseduffield/lazydocker}




\bibitem{syncatomic}
Go Authors. (2024). \textit{sync/atomic - The Go Standard Library}.
\url{https://pkg.go.dev/sync/atomic}

\bibitem{runtimememstats}
Go Authors. (2024). \textit{runtime - Memory Statistics in Go}.
\url{https://pkg.go.dev/runtime#MemStats}

\bibitem{encodingjson}
Go Authors. (2024). \textit{encoding/json - The Go Standard Library}.
\url{https://pkg.go.dev/encoding/json}




\bibitem{battleship}
Milton Bradley. (1967). \textit{Battleship (game)}.
Original pencil and paper game dates to World War I.
\url{https://en.wikipedia.org/wiki/Battleship_(game)}

\end{thebibliography}

\end{document}
